\chapter{Analog Modulation and the Classic Radio}
\label{ch:analog}

\begin{chapterintro}
For most of the twentieth century, radio meant analog modulation: a voice or a picture
impressed directly onto the amplitude, phase or frequency of a carrier. Analog systems are
still on the air---AM and FM broadcasting, aviation voice, marine and land-mobile radio,
amateur single sideband---and their ideas are everywhere inside digital radios. The
superheterodyne receiver, the phase-locked loop, the I/Q modulator, pre-emphasis, the
capture effect and the threshold effect were all discovered by engineers building analog
equipment. This chapter treats the classic systems the way their designers did: as
engineering trade-offs between power, bandwidth, receiver complexity and fidelity.
\end{chapterintro}

\section{Why modulate at all?}

A microphone produces a baseband signal $m(t)$ with energy from roughly 100\,Hz to a few
kilohertz. Radiating it directly is hopeless for three reasons. First, antennas are
efficient only when their size is comparable to a wavelength; at 1\,kHz the wavelength is
300\,km. Second, every station would occupy the same spectrum, so no two could operate in
the same region. Third, the propagation of radio waves depends strongly on frequency, and
the designer wants to choose the band that suits the service: medium waves that follow the
ground at night, VHF that stays local, microwaves that go straight to a satellite.
\term{Modulation} solves all three by translating the message to a \term{carrier} frequency
$f_c$ chosen by the engineer and allocated by the regulator.

A sinusoidal carrier $A_c\cos(2\pi f_c t+\phi)$ has only two things that can be varied:
its amplitude and its angle. Every analog scheme in this chapter varies one or both, and
every digital scheme in Part~III does the same with a discrete alphabet. In the language of
Chapter~\ref{ch:signals}, a modulated signal is
\begin{equation}
s(t)=\re\{\tilde{s}(t)\,e^{j2\pi f_c t}\},
\label{eq:ch04:env}
\end{equation}
and each modulation is nothing more than a rule for building the complex envelope
$\tilde{s}(t)$ from $m(t)$:

\begin{center}
\begin{tabular}{lll}
\toprule
Scheme & Complex envelope $\tilde s(t)$ & Bandwidth \\
\midrule
AM (DSB-LC) & $A_c[1+\mu m(t)]$ & $2W$ \\
DSB-SC & $A_c\,m(t)$ & $2W$ \\
SSB (upper) & $A_c[m(t)+j\hat m(t)]$ & $W$ \\
PM & $A_c e^{jk_p m(t)}$ & $\approx 2(\Delta f+W)$ \\
FM & $A_c e^{j2\pi k_f\int m(\tau)d\tau}$ & $\approx 2(\Delta f+W)$ \\
\bottomrule
\end{tabular}
\end{center}

Here $W$ is the message bandwidth, $\hat m(t)$ the Hilbert transform of $m(t)$ and $\Delta f$
the peak frequency deviation. The table is worth a moment's study, because it is the whole
chapter in five lines: amplitude modulations have \emph{real} or \emph{analytic} envelopes and
occupy a bandwidth tied to $W$; angle modulations have constant-magnitude envelopes and
trade bandwidth for noise immunity.

\section{Amplitude modulation}

\subsection{The AM signal}

Conventional AM, known formally as double sideband with large carrier (DSB-LC), was the
first broadcast modulation and is still used on medium and short waves and for aviation
voice around 118--137\,MHz. With the message normalised so that $|m(t)|\le 1$,
\begin{equation}
s_{\text{AM}}(t)=A_c\,[1+\mu\,m(t)]\cos(2\pi f_c t),
\end{equation}
where $0<\mu\le 1$ is the \term{modulation index}. The envelope of the RF waveform traces
$A_c[1+\mu m(t)]$, which is always positive as long as $\mu\le1$, so a diode, a capacitor and
a resistor can recover the message. Figure~\ref{fig:ch04_am} shows a tone-modulated carrier at
three indices. At $\mu=1.4$ the envelope would need to go negative; it cannot, so the detector
output folds and the audio is badly distorted. Broadcast transmitters limit negative peaks
to prevent this \term{over-modulation}, while allowing positive peaks somewhat above 100\%
because those do not break the envelope.

\mdcfig{ch04_am}{AM with a 500\,Hz tone on a 10\,kHz carrier for three modulation indices.
The dashed green curve is the ideal envelope $1+\mu m(t)$; the red curve is the output of a
simulated full-wave envelope detector. Over-modulation ($\mu=1.4$) makes the envelope fold,
and the detected audio is distorted.}

Taking the Fourier transform,
\begin{equation}
S_{\text{AM}}(f)=\frac{A_c}{2}\big[\delta(f-f_c)+\delta(f+f_c)\big]
+\frac{A_c\mu}{2}\big[M(f-f_c)+M(f+f_c)\big].
\end{equation}
The spectrum is the message spectrum copied to either side of the carrier: an
\term{upper sideband} and a \term{lower sideband}, each a complete copy of the information, plus
an unmodulated carrier line. The occupied bandwidth is $2W$; AM broadcast channels are
10\,kHz wide in the Americas and 9\,kHz elsewhere, which limits the audio to about 4.5 to
5\,kHz.

\subsection{Power efficiency}

With a tone message $m(t)=\cos 2\pi f_m t$, the carrier power is $A_c^2/2$ and each sideband
carries $A_c^2\mu^2/8$. The fraction of the transmitted power that carries information is
\begin{equation}
\eta_{\text{AM}}=\frac{\mu^2P_m}{1+\mu^2P_m}\qquad\Longrightarrow\qquad
\eta=\frac{1}{3}\ \text{for a tone at }\mu=1,
\end{equation}
where $P_m$ is the normalised message power. Speech has a high peak-to-average ratio, so with
peaks held at 100\% modulation the average $P_m$ is often below 0.1 and $\eta$ falls to a few
percent. A 50\,kW clear-channel station may radiate only a kilowatt or two of sidebands.

Why accept such waste? Because in 1920 the only practical receiver was an envelope
detector, and a large carrier makes envelope detection possible. The transmitter is
expensive but there is one of it; the receivers are cheap and there are millions. That
asymmetry---put complexity where there are few units---recurs throughout the history of
communications, from broadcast television to GPS.

\begin{inpractice}[title=Modern AM transmitters]
A modern high-power AM transmitter does not use a linear amplifier. It uses \emph{pulse
step} or \emph{digital} modulation: dozens of solid-state modules, each switched fully on or
off, whose outputs are combined so that the number of active modules follows the audio.
Switching amplifiers are 90\% efficient or better. The same idea, envelope tracking, is used
in every 4G and 5G handset to power the PA efficiently while it amplifies an OFDM signal
(Chapter~\ref{ch:sdr}).
\end{inpractice}

\subsection{The envelope detector}

The classic detector is a diode followed by a parallel $RC$ network. On each positive
half-cycle the diode charges the capacitor to the peak; between peaks the capacitor
discharges through $R$. Two constraints fix the time constant:
\begin{equation}
\frac{1}{f_c}\ll RC\ll\frac{1}{W}.
\end{equation}
If $RC$ is too short the output ripples at the carrier frequency. If it is too long the
capacitor cannot follow the falling envelope, and the output shows \emph{diagonal clipping}.
A tighter analysis for a tone at index $\mu$ gives $RC\le\sqrt{1-\mu^2}/(2\pi f_m\mu)$. For
broadcast AM with $f_c\approx 1$\,MHz and $W=5$\,kHz there is a comfortable factor of 200
between the limits.

\begin{worked}[title=Designing an envelope detector]
A 455\,kHz IF carries AM with audio up to 5\,kHz at up to 90\% modulation. The diagonal
clipping limit at 5\,kHz and $\mu=0.9$ is
$RC\le\sqrt{1-0.81}/(2\pi\cdot 5000\cdot 0.9)=15.4\,\mu$s. The ripple condition asks for
$RC\gg 1/455\,\text{kHz}=2.2\,\mu$s. Choosing $R=10$\,k$\Omega$ and $C=1$\,nF gives
$RC=10\,\mu$s, which satisfies both. The remaining 455\,kHz ripple is removed by a simple
audio low-pass filter.
\end{worked}

\subsection{Double sideband with suppressed carrier}

If the receiver can regenerate the carrier, it does not need to be transmitted. Removing
it gives DSB-SC,
\begin{equation}
s_{\text{DSB}}(t)=A_c\,m(t)\cos(2\pi f_c t),
\end{equation}
generated by a \term{balanced modulator}---a product of two signals, in practice a Gilbert
cell or a diode ring. All transmitted power is now in the sidebands, but the envelope is
$|m(t)|$, so an envelope detector no longer works. The receiver must multiply by a
\emph{coherent} local carrier $\cos(2\pi f_c t+\theta)$ and low-pass filter:
\begin{equation}
\text{LPF}\{s_{\text{DSB}}(t)\cos(2\pi f_c t+\theta)\}=\tfrac{A_c}{2}m(t)\cos\theta.
\label{eq:ch04:coh}
\end{equation}
A phase error $\theta$ scales the output by $\cos\theta$; at $90^\circ$ the output vanishes.
A frequency error $\Delta f$ makes the output beat as $m(t)\cos(2\pi\Delta f t)$. Coherent
detection therefore needs a carrier recovery loop. John Costas's 1956 answer---two
product detectors in quadrature whose outputs multiply to give a phase error---is the
\term{Costas loop} that Chapter~\ref{ch:sync} uses to recover the carrier of every BPSK and QPSK
signal.

\mdcfig{ch04_am_family}{Spectra of the four amplitude modulations with a voice-like message
on a 12\,kHz carrier: conventional AM carries a strong carrier line; DSB-SC removes it; SSB
keeps only one sideband; VSB keeps one sideband plus a gradual vestige of the other.}

\section{Single sideband}

\subsection{Only half the spectrum is needed}

For a real message, the lower sideband is the mirror image of the upper; transmitting both
wastes half the bandwidth. \term{Single sideband} (SSB) sends one. Its complex envelope is the
analytic signal of the message,
\begin{equation}
\tilde s_{\text{USB}}(t)=m(t)+j\hat m(t),\qquad
s_{\text{USB}}(t)=m(t)\cos 2\pi f_c t-\hat m(t)\sin 2\pi f_c t,
\label{eq:ch04:ssb}
\end{equation}
where $\hat m$ is the Hilbert transform (Chapter~\ref{ch:signals}). Changing the sign of the
second term gives the lower sideband. SSB occupies bandwidth $W$, puts all its power into
information, and fades more gracefully than AM because there is no carrier to suffer
selective fading. It has been the standard for long-haul HF voice since the 1950s: military,
maritime, aeronautical and amateur. Its analog-telephony cousin, SSB frequency-division
multiplexing, carried the long-distance telephone network (Chapter~\ref{ch:history}) with
12-channel groups stacked into supergroups and mastergroups on coaxial cables and microwave
relays.

\subsection{Three ways to generate SSB}

\begin{description}
\item[Filter method.] Generate DSB-SC at a low intermediate frequency and remove one
sideband with a sharp crystal or mechanical filter. The filter must pass 300\,Hz and reject
$-300$\,Hz, a transition band of 600\,Hz at, say, 9\,MHz---hence the need for crystal
lattices and the reason voice SSB traditionally cuts off below 300\,Hz.
\item[Phasing method.] Implement~\eqref{eq:ch04:ssb} directly: two balanced modulators driven
in quadrature, one fed with $m(t)$ and the other with a 90$^\circ$-shifted copy. The
unwanted sideband cancels only as well as the two paths match. Figure~\ref{fig:ch04_ssb_phase}
shows how quickly suppression degrades: a $2^\circ$ phase error limits it to about 35\,dB.
This is precisely the \emph{image rejection} problem of every I/Q transmitter and
direct-conversion receiver, including the B200 of Chapter~\ref{ch:sdr}.
\item[Weaver method.] Shift the audio band to be centred on 0\,Hz with a quadrature
oscillator at the band centre, low-pass filter each branch, then up-convert in quadrature.
It avoids both the sharp filter and the wideband 90$^\circ$ network, and it is the method most
DSP-based SSB implementations use.
\end{description}

\mdcfig[0.7]{ch04_ssb_phase}{Unwanted-sideband suppression of the phasing method as a
function of the phase error of the 90$^\circ$ network, for three gain mismatches. The same
curves give the image rejection of any I/Q modulator or zero-IF receiver.}

\subsection{Receiving SSB}

An SSB receiver is a coherent product detector with a local beat-frequency oscillator. A
frequency error $\Delta f$ does not destroy the signal as it does for DSB-SC; it shifts every
audio component by $\Delta f$. Speech stays intelligible up to about 50\,Hz of error, which is
why SSB could be used with free-running oscillators, and why a mistuned SSB voice sounds
like ``Donald Duck''. Music, whose harmonics lose their integer ratios, cannot tolerate
even a few hertz, which is one reason SSB was never adopted for broadcasting.

\section{Vestigial sideband and analog television}

Television video extends down to DC. A clean SSB filter with a transition at DC is
impossible, and DSB wastes spectrum a TV channel cannot afford. \term{Vestigial sideband}
(VSB) keeps one sideband and a gradually rolled-off \emph{vestige} of the other. If the
filter's response is antisymmetric about $f_c$---$H(f_c+f)+H(f_c-f)=$ constant over the
message band---coherent or envelope detection recovers the message without distortion.
Readers who have met the Nyquist criterion for zero intersymbol interference
(Chapter~\ref{ch:baseband}) will recognise the same symmetry condition, discovered twice.

NTSC television, used in North America from 1941 to 2009, put a VSB video carrier 1.25\,MHz
above the lower edge of a 6\,MHz channel, with 4.2\,MHz of luminance bandwidth, an FM
audio carrier 4.5\,MHz above the video carrier, and colour information on a 3.58\,MHz
subcarrier using \emph{quadrature} amplitude modulation of two colour-difference signals.
The colour subcarrier frequency was chosen as an odd multiple of half the line rate so that
its energy interleaved between the harmonics of the luminance spectrum---a frequency-domain
multiplexing trick of real elegance. The digital ATSC standard that replaced NTSC kept VSB,
now as 8-VSB carrying 19.39\,Mb/s in the same 6\,MHz.

\section{Angle modulation}

\subsection{Phase and frequency modulation}

Angle modulation keeps the amplitude constant and varies the phase:
\begin{equation}
s(t)=A_c\cos\big(2\pi f_c t+\phi(t)\big).
\end{equation}
In \term{phase modulation} (PM), $\phi(t)=k_p m(t)$. In \term{frequency modulation} (FM), the
instantaneous frequency $f_i(t)=f_c+\frac{1}{2\pi}\dot\phi(t)$ follows the message,
\begin{equation}
f_i(t)=f_c+k_f\,m(t),\qquad \phi(t)=2\pi k_f\int_{-\infty}^t m(\tau)\,d\tau.
\end{equation}
FM of $m$ is PM of the integral of $m$; the two differ only by a filter in front of the
modulator. The \term{peak frequency deviation} is $\Delta f=k_f\max|m(t)|$, and for a tone
of frequency $f_m$ the \term{modulation index} is $\beta=\Delta f/f_m$, the peak phase
excursion in radians.

\subsection{The spectrum of FM: Bessel functions}

FM is nonlinear: the spectrum of an FM signal is not a shifted copy of the message spectrum.
For a tone,
\begin{equation}
s(t)=A_c\cos\big(2\pi f_c t+\beta\sin 2\pi f_m t\big)=A_c\sum_{n=-\infty}^{\infty}J_n(\beta)\cos\big(2\pi(f_c+nf_m)t\big),
\label{eq:ch04:bessel}
\end{equation}
where $J_n$ is the Bessel function of the first kind. A single tone produces an infinite
set of sidebands spaced $f_m$ apart, with amplitudes $J_n(\beta)$
(Figure~\ref{fig:ch04_fm_bessel}). Three facts follow.
\begin{itemize}
\item The total power is constant, $\sum J_n^2(\beta)=1$: modulation redistributes power
from the carrier to the sidebands.
\item The carrier amplitude $J_0(\beta)$ vanishes at $\beta=2.405, 5.520, \dots$. Engineers use
these \emph{Bessel nulls} to calibrate deviation: raise the audio level until the carrier
disappears on a spectrum analyser, and the index is exactly 2.405.
\item For small $\beta$, $J_1(\beta)\approx\beta/2$ and higher orders are negligible:
\term{narrowband FM} looks like AM with the carrier in quadrature, and its bandwidth is $2f_m$.
\end{itemize}

\mdcfig{ch04_fm_bessel}{Top: Bessel functions $J_0$ to $J_4$ versus modulation index. Bottom:
line spectra of tone-modulated FM for four indices, including the first carrier null at
$\beta=2.405$.}

\subsection{Carson's rule}

In theory FM bandwidth is infinite; in practice sideband amplitudes fall rapidly once
$|n|>\beta+1$. John Carson's 1922 rule of thumb---ironically from a paper arguing that FM
could not reduce bandwidth, which is true---estimates the band containing about 98\% of the
power:
\begin{equation}
B_T\approx 2(\Delta f+W)=2(\beta+1)W.
\label{eq:ch04:carson}
\end{equation}
Figure~\ref{fig:ch04_carson} compares it with the exact 98\% bandwidth. FM broadcasting
uses $\Delta f=75$\,kHz and $W=15$\,kHz ($\beta=5$), so $B_T=180$\,kHz, which fits the
200\,kHz channel raster. Two-way land-mobile radio uses $\Delta f=2.5$ or 5\,kHz and $W=3$\,kHz,
fitting 12.5 or 25\,kHz channels.

\mdcfig[0.7]{ch04_carson}{Carson's rule (dashed) against the exact bandwidth containing
98\% of the power of tone-modulated FM (steps). The rule is accurate across the range of
practical modulation indices.}

\subsection{Generating FM}

The obvious FM modulator is a voltage-controlled oscillator (VCO) whose frequency follows
the message, typically an LC oscillator tuned by a varactor diode. Its centre frequency
drifts, so the VCO is enclosed in a phase-locked loop whose bandwidth is below the lowest
audio frequency: the loop holds the average frequency while the modulation passes through.
Armstrong's original \emph{indirect} method generated narrowband PM of the integrated audio
with a crystal-stable oscillator and then multiplied the frequency many times, which
multiplied the deviation by the same factor. Today most FM is generated digitally: an NCO
in an FPGA or DSP accumulates $2\pi k_f m[n]/f_s$ each sample and an I/Q modulator translates
the result to RF. A software-defined radio produces broadcast-quality stereo FM in a few
lines of code.

\subsection{Demodulating FM}

A demodulator must output a voltage proportional to instantaneous frequency. The
historical sequence of solutions is a tour of analog circuit ingenuity:
\begin{description}
\item[Slope detector] Tune a resonant circuit so the carrier sits on its skirt; frequency
variations become amplitude variations, which an envelope detector recovers.
\item[Foster--Seeley discriminator and ratio detector] Balanced versions using a
centre-tapped transformer; the ratio detector adds built-in amplitude limiting. These
dominated radios from the 1930s to the 1970s.
\item[Quadrature detector] Multiply the signal by a copy delayed through a phase-shift
network whose phase varies linearly with frequency; used in countless FM IF chips.
\item[PLL demodulator] Lock a VCO to the incoming signal; the VCO's control voltage is
the demodulated audio. PLL demodulators extend the threshold (below).
\item[Digital] With complex baseband samples $z[n]$, the phase difference between successive
samples is the instantaneous frequency:
\begin{equation}
\hat f[n]=\frac{f_s}{2\pi}\arg\big(z[n]\,z^*[n-1]\big).
\label{eq:ch04:digfm}
\end{equation}
This single line of code is the FM demodulator in GNU Radio's \texttt{Quadrature Demod} block
and in almost every SDR.
\end{description}

\section{Noise in analog systems}

\subsection{The figure of merit}

To compare modulations fairly, define the received carrier power $P_R$, the noise
density $N_0$, and the message bandwidth $W$. The reference is baseband transmission,
with output SNR $\gamma=P_R/(N_0W)$. For each scheme we compute the output SNR relative to
$\gamma$. The results for a tone message are the classic table:

\begin{center}
\begin{tabular}{lcc}
\toprule
System & Output SNR & Transmission bandwidth \\
\midrule
Baseband & $\gamma$ & $W$ \\
DSB-SC (coherent) & $\gamma$ & $2W$ \\
SSB (coherent) & $\gamma$ & $W$ \\
AM, envelope detector (above threshold) & $\dfrac{\mu^2P_m}{1+\mu^2P_m}\gamma\ \le\ \gamma/3$ & $2W$ \\
FM (above threshold) & $3\beta^2P_m\,\gamma$ & $2(\beta+1)W$ \\
\bottomrule
\end{tabular}
\end{center}

DSB-SC and SSB tie with baseband: coherent linear modulation neither gains nor loses. AM
loses at least 4.8\,dB to the carrier. FM is different in kind: its output SNR grows as
$\beta^2$, so every doubling of deviation (and roughly of bandwidth) buys 6\,dB. This is the
first appearance in this book of a principle that information theory
(Chapter~\ref{ch:infotheory}) makes exact: \emph{bandwidth can be exchanged for signal-to-noise
ratio}. For broadcast FM with $\beta=5$ and $P_m=1/2$, the improvement is $37.5$, or
15.7\,dB, before pre-emphasis.

\subsection{Why FM noise rises with frequency}

In the demodulator, additive noise perturbs the phase by an amount inversely proportional
to the carrier amplitude, and the discriminator differentiates phase. Differentiation
multiplies the noise spectrum by $f^2$, so FM's output noise PSD is \emph{parabolic}: low
audio frequencies are quiet and high frequencies are noisy. Both broadcast FM and analog
TV sound exploit this with \term{pre-emphasis}: the transmitter boosts high audio
frequencies with a first-order high-pass shelf of time constant 75\,$\mu$s (Americas, Korea)
or 50\,$\mu$s (Europe and most of the world), and the receiver applies the matching
de-emphasis. Because speech and music have little energy at high frequencies, the boost
hardly increases the deviation, while de-emphasis cuts the parabolic noise by a further
10 to 13\,dB. Dolby noise reduction in cassette recorders used the same idea.

\subsection{The threshold effect}

The FM advantage holds only above a carrier-to-noise ratio of about 10\,dB in the IF
bandwidth. Below it, the noise occasionally drives the received phasor around the origin,
adding a full $2\pi$ to the phase. The discriminator turns each of these \term{clicks} into
an impulse, and the output SNR collapses. Figure~\ref{fig:ch04_fm_threshold} shows a Monte
Carlo simulation: above threshold, the curves follow the $3\beta^2P_m\gamma$ theory and higher
deviation wins; below it, the curves bend sharply down, and the wider-band signal crosses
threshold first because it collects more noise. The threshold is the analog ancestor of the
``cliff effect'' of digital systems: a good FM link sounds perfect until it suddenly does
not. PLL and feedback demodulators lower the threshold by a few decibels by effectively
narrowing the noise bandwidth; they were developed for the space programme, where every
decibel of a deep-space link mattered.

\mdcfig{ch04_fm_threshold}{Output SNR versus carrier-to-noise ratio for FM with
$\beta=2$ and $\beta=5$ (Monte Carlo with a tone), compared with AM and baseband. Above
threshold FM beats AM by up to 20\,dB; below about 10\,dB CNR the ``clicks'' take over.}

\subsection{The capture effect}

When two FM signals arrive on the same channel, the demodulator output is dominated by the
stronger one, and the weaker is suppressed almost completely once the ratio exceeds a few
decibels (the \emph{capture ratio}, 1 to 2\,dB in a good receiver). AM has no such effect: an
interferer 20\,dB weaker is still clearly heard. Capture is why FM made VHF mobile radio
practical: co-channel stations could be reused much closer together. It is also why aircraft
still use AM on VHF: two pilots transmitting simultaneously produce an audible heterodyne on
AM, whereas FM would silently hide the weaker transmission.

\section{FM broadcasting in detail}

\subsection{Stereo multiplex}

FM stereo, standardised in the USA in 1961, had to remain compatible with millions of mono
receivers. The solution (Figure~\ref{fig:ch04_stereo_mpx}) builds a \term{multiplex} (MPX)
baseband that frequency-modulates the carrier:
\begin{itemize}
\item the sum signal $L+R$ from 30\,Hz to 15\,kHz, which a mono receiver plays as usual;
\item a 19\,kHz \term{pilot} tone at 8 to 10\% of deviation;
\item the difference $L-R$ as DSB-SC on a suppressed 38\,kHz subcarrier, occupying 23 to
53\,kHz.
\end{itemize}
The stereo decoder doubles the pilot to regenerate the 38\,kHz carrier coherently, recovers
$L-R$, and forms $L=\frac12[(L+R)+(L-R)]$ and $R=\frac12[(L+R)-(L-R)]$. Because the $L-R$ channel
sits high in the MPX spectrum where FM noise is parabolic, stereo reception is about 20\,dB
noisier than mono, which is why receivers blend to mono on weak signals.

Above the stereo band, the \term{Radio Data System} (RDS, or RBDS in North America) carries
station names, programme type and traffic information as differentially encoded BPSK at
1187.5\,b/s on a 57\,kHz subcarrier, the third harmonic of the pilot. Older subsidiary
carriers at 67 and 92\,kHz carried background music and reading services. The MPX spectrum is
a compact lesson in frequency-division multiplexing, coherent carrier regeneration, and
backward compatibility.

\mdcfig{ch04_stereo_mpx}{The FM stereo multiplex baseband: mono-compatible $L+R$, the 19\,kHz
pilot, the $L-R$ DSB-SC signal around a suppressed 38\,kHz subcarrier, and RDS at 57\,kHz.}

\subsection{HD Radio and DAB}

Digital radio arrived in two philosophies. Europe's DAB (1995) and its successor DAB+ use
OFDM in new VHF spectrum, multiplexing a dozen services in a 1.5\,MHz ensemble with
single-frequency networks. The US HD Radio system instead places low-power OFDM sidebands on
either side of existing AM and FM analog carriers (``in-band on-channel''), allowing a
gradual transition without new spectrum. Both use the OFDM and coding machinery of
Part~V.

\section{The superheterodyne receiver}

\subsection{Armstrong's architecture}

Early receivers amplified directly at the signal frequency in several tuned stages
(the tuned radio-frequency, or TRF, receiver). Tracking several tuned circuits across a
band was hard, and selectivity and gain varied with tuning. Edwin Armstrong's
\term{superheterodyne}, conceived in France in 1918, mixes every incoming station down to a
single fixed \term{intermediate frequency} (IF), where all the selectivity and most of the
gain can be designed once and for all. Nearly every radio receiver built since 1930 is a
superheterodyne or a direct descendant of one.

A mixer multiplies the input with a local oscillator (LO). An input at $f_{RF}$ and an LO at
$f_{LO}$ produce $|f_{RF}\pm f_{LO}|$; the IF filter selects $f_{IF}=|f_{RF}-f_{LO}|$. With
\emph{high-side injection}, $f_{LO}=f_{RF}+f_{IF}$, so tuning the receiver means tuning the LO.

\subsection{The image problem}

The mixer cannot tell whether the input is above or below the LO: an input at
$f_{\text{image}}=f_{LO}+f_{IF}$, two IFs away from the wanted signal, produces the same
IF. Figure~\ref{fig:ch04_superhet} illustrates an FM receiver tuned to 98.1\,MHz with a
10.7\,MHz IF: the image is at 119.5\,MHz, which falls in the aeronautical band. A tuned
\term{preselector} ahead of the mixer must attenuate the image before mixing; once mixed, the
image is indistinguishable from the wanted station. This sets the first rule of IF
selection: \emph{the IF must be high enough that the preselector can reject the image}.
The second rule is its opposite: \emph{the IF must be low enough that the IF filter can
provide adjacent-channel selectivity}. At 455\,kHz, a ceramic filter easily makes a 10\,kHz
AM channel; at 10.7\,MHz, a ceramic filter makes a 200\,kHz FM channel.

\mdcfig{ch04_superhet}{Frequency plan of an FM superheterodyne with a 10.7\,MHz IF and
high-side LO. The preselector must suppress the image two IFs above the wanted station before
the mixer.}

\subsection{Double conversion and the path to zero IF}

When one IF cannot satisfy both rules, use two. Communications receivers convert first to
a high IF (often above the tuning range, so the image is far away and a fixed low-pass
filter rejects it) and then to a low IF where the channel filtering is done. Cellular
receivers of the 1990s followed this pattern. The logical end point is to convert directly
to an IF of zero: the \term{direct-conversion} or zero-IF receiver, in which the ``image'' is
the signal's own negative-frequency half and is rejected not by a filter but by I/Q
balance. That architecture, with its DC-offset and image-rejection problems, is the subject
of Chapter~\ref{ch:sdr}.

\subsection{Automatic gain control}

Received signal strengths span more than 100\,dB, from a station across town to one at the
edge of coverage, and fading adds rapid swings of tens of decibels. \term{Automatic gain
control} (AGC) measures the IF level and adjusts the gain of the RF and IF stages to keep the
detector input constant. AGC time constants encode a compromise: fast enough to follow
fading, slow enough not to flatten the speech envelope of AM. In digital receivers the same
loop protects the ADC from clipping and is the first block in every receive chain.

\section{Analog systems that survive}

Analog modulation is not only history. In 2026 it remains in service where simplicity,
backward compatibility or graceful degradation matter:
\begin{itemize}
\item AM broadcasting on long, medium and short waves, and aviation VHF voice (AM), for the
capture-effect reason above and for decades of installed equipment.
\item FM broadcasting worldwide, although some countries (Norway in 2017, Switzerland by
2026) have switched national FM networks off in favour of DAB+.
\item Narrowband FM in marine VHF and amateur radio, although professional land-mobile radio
has largely moved to digital systems such as P25, TETRA and DMR.
\item SSB on HF for maritime, military and amateur long-distance links, increasingly
alongside digital HF modems.
\item Analog FM links inside satellite transponder tests, in some wireless microphones and in
radio astronomy calibration.
\end{itemize}

\begin{keyidea}[title=What analog radio teaches the digital engineer]
\begin{itemize}
\item \textbf{Everything is a complex envelope.} AM, SSB, PM and FM differ only in how
$\tilde{s}(t)$ is built. So do QPSK, QAM and OFDM.
\item \textbf{Coherent detection needs a carrier reference.} The Costas loop was invented
for DSB-SC voice.
\item \textbf{Bandwidth buys noise immunity}, but only above a threshold. Shannon made both
statements exact.
\item \textbf{Image rejection and I/Q balance are the same problem}, from the phasing SSB
exciter of 1950 to the zero-IF RFIC of today.
\end{itemize}
\end{keyidea}

\begin{labbox}[title=Lab: analog modulation in Python]
The notebook \lab{labs/ch04\_analog\_modulation.ipynb} builds AM, DSB-SC, SSB (phasing and
Weaver) and FM modulators and demodulators on complex baseband samples; lets you vary
modulation index and noise with sliders while listening-free metrics (SINAD, THD) update;
reproduces the FM threshold curve; and decodes a synthetic FM stereo MPX signal with
pilot-based 38\,kHz regeneration. With the B200, \lab{gnuradio/gr01\_spectrum\_iq\_capture.py}
records a real broadcast FM station that the notebook can demodulate.
\end{labbox}

\problems
\begin{enumerate}[label=\thechapter.\arabic*]
\item An AM transmitter radiates 10\,kW of carrier. With tone modulation at $\mu=0.8$, what is
the total power and the power in each sideband? What is the efficiency?
\item Show that the envelope detector's diagonal-clipping condition for a tone is
$RC\le\sqrt{1-\mu^2}/(2\pi f_m\mu)$.
\item A DSB-SC signal is demodulated with a local carrier having a phase error that drifts
linearly at 1\,Hz. Describe the output. Repeat for SSB.
\item Derive the Bessel expansion~\eqref{eq:ch04:bessel} from the generating function
$e^{j\beta\sin\theta}=\sum_n J_n(\beta)e^{jn\theta}$.
\item Using Carson's rule, how many 12.5\,kHz land-mobile FM channels ($\Delta f=2.5$\,kHz, $W=3$\,kHz)
fit in the bandwidth of one broadcast FM channel? Compare the output SNR advantage over AM
of each system.
\item An FM receiver has IF 10.7\,MHz and high-side LO. It is tuned to 88.1\,MHz. Find the
image frequency. What preselector attenuation is needed if an image signal is 40\,dB stronger
than the wanted one and the demodulator needs 20\,dB of protection?
\item Show that the digital FM demodulator~\eqref{eq:ch04:digfm} is unbiased for any carrier
amplitude and explain what happens when $|\hat f|>f_s/2$.
\item Explain why FM stereo is about 20\,dB noisier than mono at the same carrier-to-noise
ratio. (\emph{Hint:} integrate the parabolic noise PSD over 23--53\,kHz and over 0--15\,kHz.)
\item \emph{(Simulation)} Generate an FM stereo MPX signal with different tones in $L$ and $R$,
frequency-modulate it with 75\,kHz deviation at 400\,kS/s, add noise, and write a decoder.
Plot left-right separation versus pilot phase error.
\end{enumerate}

\furtherreading
\begin{itemize}
\item B.~P.~Lathi and Z.~Ding, \emph{Modern Digital and Analog Communication Systems}, 5th ed.,
Oxford, 2019, Ch.~4--5. A clear textbook treatment of AM and FM with noise analysis.
\item S.~Haykin and M.~Moher, \emph{Communication Systems}, 5th ed., Wiley, 2009.
\item E.~H.~Armstrong, ``A method of reducing disturbances in radio signaling by a system of
frequency modulation,'' \emph{Proc. IRE}, 24(5), 1936. The paper that launched FM.
\item J.~R.~Carson, ``Notes on the theory of modulation,'' \emph{Proc. IRE}, 10, 1922.
\item U.~L.~Rohde, J.~C.~Whitaker and H.~Zahnd, \emph{Communications Receivers: Principles and
Design}, 4th ed., McGraw-Hill, 2017. The practitioner's reference on superheterodynes.
\item \emph{The ARRL Handbook for Radio Communications} (annual). Practical SSB, FM and receiver
circuits, updated every year.
\end{itemize}
